\documentclass[11pt,a4paper]{article}
\usepackage[margin=2.4cm]{geometry}
\usepackage{amsmath,amsthm,mathtools}
\usepackage{fontspec}
\usepackage{unicode-math}
\setmainfont{STIX2Text}[Path=fonts/, Extension=.otf,
  UprightFont=*-Regular, ItalicFont=*-Italic,
  BoldFont=*-Bold, BoldItalicFont=*-BoldItalic]
\setmathfont{STIX2Math}[Path=fonts/, Extension=.otf]
\usepackage[protrusion=true,expansion=false]{microtype}
\usepackage{booktabs,enumitem}
\usepackage{placeins}
\usepackage{tikz}
\usetikzlibrary{arrows.meta,positioning,calc,fit,backgrounds}
\definecolor{tierspend}{RGB}{200,60,60}
\definecolor{tierfull}{RGB}{210,140,40}
\definecolor{tierin}{RGB}{60,130,180}
\definecolor{tierrecv}{RGB}{70,150,90}
\usepackage[psdextra,hidelinks]{hyperref}

\emergencystretch=3em
\hbadness=10000
\hfuzz=2pt

\newtheorem{theorem}{Theorem}[section]
\newtheorem{lemma}[theorem]{Lemma}
\newtheorem{proposition}[theorem]{Proposition}
\newtheorem{corollary}[theorem]{Corollary}
\theoremstyle{definition}
\newtheorem{definition}[theorem]{Definition}
\newtheorem{construction}[theorem]{Construction}
\newtheorem{assumption}[theorem]{Assumption}
\newtheorem{remark}[theorem]{Remark}
\newtheorem{example}[theorem]{Example}

% kind-coded theorem blocks, palette shared with the figures and the website:
% definitions/constructions blue (tierin), assumptions amber (tierfull),
% proved statements a neutral bar; remarks and examples run unframed
\usepackage{mdframed}
\providecolor{tierin}{RGB}{60,130,180}
\providecolor{tierfull}{RGB}{210,140,40}
\mdfdefinestyle{kindbar}{topline=false,bottomline=false,rightline=false,
  linewidth=1.5pt,innerleftmargin=9pt,innerrightmargin=4pt,
  innertopmargin=5pt,innerbottommargin=6pt,leftmargin=0pt,rightmargin=0pt,
  skipabove=10pt,skipbelow=10pt}
\surroundwithmdframed[style=kindbar,linecolor=tierin!70,backgroundcolor=tierin!4]{definition}
\surroundwithmdframed[style=kindbar,linecolor=tierin!70,backgroundcolor=tierin!4]{construction}
\surroundwithmdframed[style=kindbar,linecolor=tierfull!80,backgroundcolor=tierfull!5]{assumption}
\surroundwithmdframed[style=kindbar,linecolor=black!30]{theorem}
\surroundwithmdframed[style=kindbar,linecolor=black!30]{lemma}
\surroundwithmdframed[style=kindbar,linecolor=black!30]{proposition}
\surroundwithmdframed[style=kindbar,linecolor=black!30]{corollary}

\usepackage{fancyhdr}
\pagestyle{fancy}
\fancyhf{}
\fancyhead[L]{\footnotesize\scshape The Zcash Arboretum}
\fancyhead[R]{\footnotesize\itshape\nouppercase{\leftmark}}
\fancyfoot[C]{\footnotesize\thepage}
\renewcommand{\headrulewidth}{0.3pt}
\renewcommand{\sectionmark}[1]{\markboth{\thesection\ \ #1}{}}

\newcommand{\F}{\mathbb{F}}
\newcommand{\Z}{\mathbb{Z}}
\newcommand{\N}{\mathbb{N}}
\newcommand{\G}{\mathbb{G}}
\newcommand{\E}{\mathbb{E}}
\newcommand{\PP}{\mathbb{P}}
\renewcommand{\Pr}{\mathrm{Pr}}
\newcommand{\Fp}{\F_p}
\newcommand{\Fq}{\F_q}
\newcommand{\ip}[2]{\langle #1,\, #2 \rangle}
\newcommand{\Comm}{\mathsf{Commit}}
\newcommand{\Open}{\mathsf{Open}}
\newcommand{\Setup}{\mathsf{Setup}}
\newcommand{\Verify}{\mathsf{Verify}}
\newcommand{\negl}{\mathsf{negl}}
\newcommand{\poly}{\mathsf{poly}}
\newcommand{\PPT}{\mathsf{PPT}}
\newcommand{\Adv}{\mathcal{A}}
\newcommand{\Prov}{\mathcal{P}}
\newcommand{\Ver}{\mathcal{V}}
\newcommand{\Sim}{\mathcal{S}}
\newcommand{\Ext}{\mathcal{E}}
\newcommand{\Rel}{\mathcal{R}}
\newcommand{\Lang}{\mathcal{L}}
\newcommand{\nfsym}{\mathsf{nf}}
\newcommand{\cmsym}{\mathsf{cm}}
\newcommand{\cmx}{\mathsf{cmx}}
\newcommand{\cvsym}{\mathsf{cv}}
\newcommand{\rk}{\mathsf{rk}}
\newcommand{\Pallas}{\ensuremath{\mathsf{Pallas}}}
\newcommand{\Vesta}{\ensuremath{\mathsf{Vesta}}}
\newcommand{\Ep}{\ensuremath{\mathcal{E}}}
\newcommand{\NoteCommit}{\ensuremath{\mathsf{NoteCommit}}}
\newcommand{\Sinsemilla}{\ensuremath{\mathsf{Sinsemilla}}}
\newcommand{\ShortCommit}{\ensuremath{\mathsf{ShortCommit}}}
\newcommand{\CommitIvk}{\ensuremath{\mathsf{Commit}^{\mathsf{ivk}}}}
\newcommand{\Extract}{\ensuremath{\mathsf{Extract}}}
\newcommand{\Acc}{\ensuremath{\mathsf{Acc}}}
\newcommand{\GroupHash}{\ensuremath{\mathsf{GroupHash}}}
\newcommand{\ask}{\ensuremath{\mathsf{ask}}}
\newcommand{\aksym}{\ensuremath{\mathsf{ak}}}
\newcommand{\nksym}{\ensuremath{\mathsf{nk}}}
\newcommand{\ivk}{\ensuremath{\mathsf{ivk}}}
\newcommand{\fvk}{\ensuremath{\mathsf{fvk}}}
\newcommand{\ovk}{\ensuremath{\mathsf{ovk}}}
\newcommand{\dk}{\ensuremath{\mathsf{dk}}}
\newcommand{\pkd}{\ensuremath{\mathsf{pk_d}}}
\newcommand{\gd}{\ensuremath{\mathsf{g_d}}}
\newcommand{\rcv}{\ensuremath{\mathsf{rcv}}}
\newcommand{\rcm}{\ensuremath{\mathsf{rcm}}}
\newcommand{\rivk}{\ensuremath{\mathsf{r_{ivk}}}}
\newcommand{\rsk}{\ensuremath{\mathsf{rsk}}}

\renewenvironment{abstract}{%
  \small\begin{center}{\bfseries\abstractname}\end{center}\medskip\noindent\ignorespaces
}{\par\medskip}

\title{\textbf{\Huge FlyClient Guide}\\[6pt]\large Sampling an authenticated proof-of-work history}
\author{m@rek.onl}
\date{}
\begin{document}
\maketitle
\thispagestyle{empty}
\begin{abstract}
This volume of \emph{The Zcash Arboretum} constructs the sampling argument
behind FlyClient. It assumes probability from the \emph{Math Guide},
authenticated histories from the \emph{Crypto Guide}, and work, chain
selection and block-history commitments from the \emph{Consensus Guide}.
The endpoint is a verifier that checks a committed candidate history by
sampling, together with an exact statement of what would still be needed
to instantiate that verifier for Zcash. The sampling theorem is proved
here. The mining assumption, variable-difficulty transition argument and
network availability premise are kept separate from it.
\end{abstract}
\clearpage
\tableofcontents
\clearpage

\section{The chain-verification contract}\label{sec:fc-problem}

The probability language used here---random choices, events,
independence, and error bounds---is constructed in the \emph{Math Guide},
\S\emph{Probability and computation}. That section also constructs real
intervals, integration, probability densities, logarithms, and inverse
sampling before they are used below. \emph{Crypto Guide},
\S\emph{Security contracts and reductions}, supplies the computational
meaning of a negligible failure probability and a security assumption.

A full node checks blocks and their transactions before using their work
to select a chain. A resource-limited client would like to learn which
history this selects without receiving every header. Receiving a short
history commitment solves only half that problem: a hash binds a server
to its answer, but does not make the answer valid.

FlyClient adds unpredictable tests of the committed history. A server
first fixes a candidate head. The client then asks for selected earlier
headers and checks both their work and their connection to that head.
A candidate containing enough invalid work should fail a test. The
distribution must catch forks near the tip as well as forks near genesis.

Our source is B\"unz, Kiffer, Luu and Zamani,
\href{https://eprint.iacr.org/2019/226}{\emph{FlyClient: Super-Light Clients
for Cryptocurrencies}}, revision 2020-08-20, especially Sections 4--6
and Appendix A. We use zero-based history indices consistently and
separate the paper's abstract construction from the Zcash history tree.
The deployed tree's exact encoding belongs to the \emph{Consensus Guide},
\S\emph{Authenticated block history}; repeating it here would obscure the
different object the sampling theorem assumes.

\subsection{Three different questions}

Let a candidate consist of headers
\[
 C=(B_0,\ldots,B_n).
\]
The head is \(B_n\). The preceding \(n\) headers form its committed
history. A header's work is a positive integer \(w_i\); work and target
are not synonyms. The \emph{Consensus Guide},
\S\emph{Work, difficulty, and chain selection}, constructs Zcash's work
function and explains why a larger numerical target means an easier
proof of work.

The client must distinguish three predicates, meaning conditions that
can be true or false:
\begin{enumerate}
\item \emph{Membership:} a particular header is at a particular position
  in the sequence committed by \(B_n\).
\item \emph{Header validity:} the candidate obeys its header-link,
  proof-of-work and difficulty rules.
\item \emph{Ledger validity:} transactions, proofs, signatures, balances,
  nullifiers and all other state transitions are valid.
\end{enumerate}
An authenticated path proves the first predicate. Sampling aims at the
second, under explicit mining and commitment assumptions. Neither
operation independently checks every transaction in the third.

A client accepting a transaction-inclusion proof therefore relies on the
underlying consensus model's honest full nodes to reject invalid ledger
transitions. The transaction's inclusion path does not itself prove that
its spend was authorised. This distinction is the same one that separates
a transaction digest from a valid transaction.

\begin{definition}[Client outcome]\label{def:fc-outcome}
The client receives candidate heads and their proofs from servers, at
least one of which supplies an honest full node's chain. It rejects
malformed candidates and selects a surviving candidate of greatest
verified claimed work. Security asks that the selected history agree
with the honest chain away from its potentially unstable suffix.
\end{definition}

The availability premise is substantive. A completely isolated client
can receive a valid but old history and no evidence that a newer one
exists. No sampling distribution manufactures a missing honest view.
Nor does agreement between two server responses prove that either server
is independent or honest. The client checks the claimed objects; it does
not replace validation by counting matching responses.

\subsection{The assumption paid for by sampling}

Suppose a competing history branches from the honest history at some
point \(a\). A server claiming at least the honest chain's work cannot
ordinarily manufacture all of its fork with valid puzzles. FlyClient
formalises the resulting deficit through a bound on the fraction of
post-fork work that can pass the sample verifier.

\begin{assumption}[Bounded valid fork weight]\label{ass:fc-fork}
Fix \(0<c<1\) and a positive integer short-fork cutoff \(L\).
Except with probability
\(\varepsilon_{\mathrm{mine}}\), every adversarial competing fork outside
the short-fork exception has at most a \(c\) fraction of its claimed
post-fork work in headers that pass the protocol's work, history and
difficulty checks. The remainder is detectable invalid work.
\end{assumption}

This is a property of adversarial histories, not a definition of the
adversary's hash rate. In particular, \(c=0.9\) does not mean that
90 percent of total mining power is adversarial. A theorem connecting
hash-rate and network parameters to \(c,L\) must be established in the
chosen chain model.

Short forks are excluded because mining has variance even under honest
majority. The verifier checks a trailing region in full. In a
constant-difficulty chain its length can be compared directly with \(L\).
In a variable-difficulty chain, a block-count cutoff and a work-fraction
cutoff are different quantities. A concrete instantiation must make its
fully checked suffix cover the short forks exempted by its mining
assumption.

\section{Work intervals and authenticated samples}\label{sec:fc-diffmmr}

The \emph{Crypto Guide}, \S\emph{Authenticated trees and append-only
histories}, supplies ordered inclusion proofs and append-only
commitments. We now specify what their interface must expose to the
sampler. This is additional metadata on the authenticated sequence, not
a new security assumption about the hash function.

\subsection{From positions to work}

Define prefix work and total history work by
\[
 W_0=0,\qquad W_i=\sum_{j=0}^{i-1}w_j,\qquad W=W_n.
\]
For \(W>0\), header \(B_i\) occupies the half-open interval
\[
 I_i=[W_i/W,W_{i+1}/W).
\]
These intervals are disjoint and cover \([0,1)\). A point \(x\) selects
the unique header whose interval contains it. Half-open intervals fix
boundary behaviour without double counting.

For equal work per block, \(I_i=[i/n,(i+1)/n)\), so work sampling reduces
to position sampling. For unequal work, a high-work block covers more
of the domain. Sampling header positions uniformly would allow a server
to hide its claimed work in a small number of artificially difficult
headers.

\begin{construction}[Authenticated weight lookup]\label{con:fc-lookup}
Each tree node commits to its ordered children and their total weights.
To locate a work coordinate \(u=xW\), start with offset zero at the root.
At a parent with left weight \(W_\ell\), descend left when
\(u<W_\ell\). Otherwise subtract \(W_\ell\), add \(W_\ell\) to the
offset, and descend right. The proof supplies the sibling node at every
step. The verifier recomputes each parent hash and weight, then checks
that the final residual coordinate lies in the returned leaf's interval.
\end{construction}

The verifier checks lengths, child orientation, integer widths and
overflow conditions before using supplied weights. Its final total must
equal the total authenticated by the selected head. A server cannot
choose a different leaf after seeing the sample without either changing
the root or presenting an inconsistent path.

This claim is conditional on the node invariants. A hash commits to
false metadata just as effectively as true metadata. Recomputing the
visited ancestors checks the additions along that path; it does not
recompute every unopened subtree. The chain-validation argument must
account for malformed subtrees through the sample predicate.

\subsection{A history inside every header}

The abstract protocol requires header \(B_i\) to commit to the ordered
prefix \((B_0,\ldots,B_{i-1})\), not merely to its predecessor's hash.
Let that commitment be \(R_i\). The head's history root is \(R_n\).

A sample proof has two jobs. First it proves that \(B_i\) appears under
\(R_n\). Second it reconstructs the commitment to the prefix strictly
before \(B_i\), and checks that this equals \(R_i\) inside the sampled
header. Without the second check, a server could try to insert a
previously mined header into a different committed history.

For an append-only forest, a prefix-recovery proof supplies the earlier
forest peaks and the left siblings inside the peak containing leaf
\(i\). Together these partition the leaves before \(i\) into complete
subtrees. Their authentication must connect them to the head's root.
Ordering these subtrees from left to right and applying the forest's
fixed root-folding rule reconstructs the prefix root. The route and
leaf count determine their positions; supplied direction bits cannot
redefine the leaf position.

The paper uses its recursively rooted MMR and the corresponding
prefix-recovery procedure. Zcash uses the different bagging convention
constructed in \emph{Consensus Guide}, \S\emph{Authenticated block history}.
The generic requirement is the same, but a proof for one root convention
is not an encoding of a proof for the other.

\begin{proposition}[History consistency of a sampled header]
\label{prop:fc-prefix}
Assume the ordered history commitment is binding. If a sampled header
passes inclusion under \(R_n\) and its own \(R_i\) equals the recovered
prefix root, then its committed earlier history is the same ordered
prefix as that authenticated by the head, except with the commitment's
binding failure probability.
\end{proposition}
\begin{proof}
Two different ordered prefixes giving the recovered root would violate
binding. The equality with \(R_i\) therefore identifies one prefix in
both commitments. Position binding additionally prevents interpreting
the same opening as a different prefix length.
\end{proof}

The proposition says nothing about the transactions inside that prefix.
It prevents reuse of a header in an incompatible history; it does not
turn the header into a proof of the entire preceding computation.

\subsection{Difficulty transitions}

Consider the paper's epoch model. An epoch contains \(M\) blocks.
A target \(T\) is replaced after elapsed time \(\Delta\) by
\[
 T'=\operatorname{clamp}\left(T\Delta f/M,\ T/\tau,\ \tau T\right),
 \qquad \tau>1,
\]
where a round is the model's unit of time and \(f\) is the intended
blocks-per-round rate. Here
\(\operatorname{clamp}(x,a,b)=\max(a,\min(x,b))\).
Rounding and the
permitted target range are part of a concrete rule. The clamp bounds
each target change; it is essential to the model's security analysis.

The paper's difficulty tree aggregates a tuple
\[
 (h,w,t,D_{\mathrm{start}},D_{\mathrm{next}},n).
\]
Here \(h\) authenticates the node, \(w\) is total difficulty weight,
\(t\) is elapsed time, \(n\) is leaf count, and \(D\) denotes difficulty
rather than the oppositely ordered proof-of-work target. A leaf has
\(n=1\), its own weight and elapsed time, and the next difficulty
prescribed by the rule. Joining left and right nodes adds \(w,t,n\),
inherits the left starting difficulty and right next difficulty, and
authenticates both full tuples.

At an opened join, the verifier checks
\(D_{\mathrm{next},\ell}=D_{\mathrm{start},r}\). Within one epoch it
checks constant difficulty and consistent aggregate weight. Across
epochs it checks that the boundary difficulties, elapsed time and total
weight admit legal clamped transitions.

For example, if a sequence of \(r\) epoch transitions begins at
difficulty \(D\), each transition changes difficulty by at most a factor
\(\tau\). After \(j\) transitions its difficulty lies between
\(D\tau^{-j}\) and \(D\tau^j\). An independently fixed ending difficulty
adds the same constraints backwards. The maximum possible intermediate
difficulty is bounded by the smaller of the forward and backward
upper envelopes; the minimum is bounded by the larger lower envelope.
Summing these envelopes bounds aggregate weight. Timing constraints
must also be satisfied: a sequence of allowable difficulty levels alone
does not show that the retarget formula produces them.

This explains the purpose of the metadata, but is not a complete
validator for arbitrary retarget algorithms. The paper's
variable-difficulty argument imports a transition-rewriting lemma:
an adversary whose malformed transitions escape some local tests can
be replaced, without materially increasing its mining budget, by one
respecting the model's transitions and with the same sample-valid work.
The lemma's hypotheses include a sufficient aggregate-feasibility test.
Our sampling theorem below assumes the resulting valid-weight bound;
it does not independently prove that lemma, and does not instantiate
its feasibility test for Zcash's averaging-window retarget.

\section{Choosing the sampling distribution}\label{sec:fc-sampling}

The real-number and continuous-probability construction in \emph{Math Guide},
\S\emph{Probability and computation}, now gives an exact mathematical
problem. Work is normalised to
\([0,1]\), and the final fraction \(\delta\) is checked in full. The
remaining samples are independent draws from a density on
\([0,1-\delta]\). A density \(g\) assigns an interval \([a,b]\) probability
\(\int_a^b g(x)\,dx\); the integral is its area. The probability rules
used here are those of \emph{Math Guide},
\S\emph{Probability and computation}.

\subsection{Why the tip receives more samples}

A fork at \(a\) occupies weight \(1-a\). Under
Assumption~\ref{ass:fc-fork}, at least
\((1-c)(1-a)\) of this weight is invalid. Uniform sampling over the whole
history catches it with probability only
\((1-c)(1-a)\), which tends to zero as the fork moves towards the tip.

Choose an increasing density. For a fixed amount of invalid work after
\(a\), its least detectable placement is at the beginning of the fork,
where the density is smallest. Moving any invalid mass from a later,
higher-density region to an available earlier region cannot increase
the catch probability. Repeating such moves leaves the interval
\[
 J_a=[a,\ 1-c(1-a)].
\]
Its length is exactly \((1-c)(1-a)\). Extra invalid work only helps the
verifier.

If this interval cannot fit before \(1-\delta\), then any placement
avoiding the fully checked suffix is impossible. Otherwise the desired
density should assign the same probability to \(J_a\) for every possible
fork point.

\subsection{Normalisation and the inverse sampler}

Use the natural logarithm constructed in \emph{Math Guide},
\S\emph{Probability and computation}: for \(y>0\),
\(\ln y=\int_1^y dt/t\). That construction proves the substitution
rule, \(\ln(ab)=\ln a+\ln b\), and the inverse exponential function.
In particular, \(\ln y\) is increasing, is zero at \(1\), and is
negative for \(0<y<1\).

Take \(0<\delta<1\) and define
\[
 g_\delta(x)=
 \frac{1}{(1-x)\ln(1/\delta)}
 \quad(0\leq x\leq1-\delta).
\]
The denominator is positive. Substituting \(t=1-x\) gives
\[
 \int_0^z g_\delta(x)\,dx
 =\frac{-\ln(1-z)}{\ln(1/\delta)}=:G_\delta(z),
 \qquad
 G_\delta(1-\delta)=1.
\]
Thus \(g_\delta\) is a probability density. It increases towards the
tip, but its support stops before the singularity at \(x=1\).

To draw a point, draw an independent uniform \(U\in[0,1]\) and solve
\(G_\delta(x)=U\):
\[
 x=1-\delta^U.
\]
Indeed, \(G_\delta\) is increasing, so
\(\Pr[1-\delta^U\leq z]=G_\delta(z)\). A finite implementation must
specify how its random bits approximate this ideal draw; floating-point
rounding is not part of the mathematical construction.

For a block occupying \([a,b)\), its probability is the integral over
that interval intersected with the sampling support. When
\(0\leq a<b\leq1-\delta\), this is
\[
 p_{a,b}=
 \frac{\ln((1-a)/(1-b))}{\ln(1/\delta)}.
\]
If the boundary cuts a block, fully checking that entire block is safe.
It enlarges the fully checked region, never reducing its protection.

\begin{theorem}[Single-sample detection]\label{thm:fc-detection}
Suppose a fork after \(a\) contains at least
\((1-c)(1-a)\) invalid work. If no invalid work occurs in the fully checked
suffix, a sample from \(g_\delta\) catches invalid work with probability
at least
\[
 p=\frac{\ln(1/c)}{\ln(1/\delta)}.
\]
The claim applies when \(\delta\leq c\). If the invalid work cannot fit
before the checked suffix, the suffix itself detects the fork.
\end{theorem}
\begin{proof}
By the increasing-density argument, the least detectable admissible
placement is \(J_a\). If it lies in the support, its mass is
\[
 \int_a^{1-c(1-a)}
 \frac{dx}{(1-x)\ln(1/\delta)}
 =
 \frac{\ln((1-a)/(c(1-a)))}{\ln(1/\delta)}
 =p.
\]
If its right endpoint exceeds \(1-\delta\), there is insufficient
uninspected work after \(a\) to contain all invalid work. Any admissible
placement then intersects the suffix.
\end{proof}

The proof works for whole blocks as well. Their invalid-work regions are
finite unions of intervals, a restricted family of the placements
allowed to the adversary in the fractional-work argument.

\subsection{Optimality and its exact domain}

The following optimality question asks for the best guarantee the
verifier can achieve against the adversary's worst permitted placement.
This best-against-worst criterion is called \emph{minimax}.

\begin{theorem}[Minimax optimality]\label{thm:fc-optimal}
Let \(k\) be a positive integer and \(\delta=c^k\).
Among all one-sample densities supported on \([0,1-\delta]\),
the largest guaranteed detection probability against the continuum
fork model is \(1/k\). The density \(g_\delta\) attains it.
\end{theorem}
\begin{proof}
For \(i=0,\ldots,k-1\), consider the fork point \(a_i=1-c^i\).
An admissible adversary makes exactly the interval
\[
 [1-c^i,1-c^{i+1}]
\]
invalid and places its valid fraction afterwards. These \(k\) intervals
partition the sampling support, apart from endpoints of zero probability
under a density. At least one interval receives probability at most
\(1/k\). Thus no density guarantees more than \(1/k\).
Theorem~\ref{thm:fc-detection} gives
\(\ln(1/c)/\ln(1/c^k)=1/k\) for \(g_\delta\).
\end{proof}

The integer condition is used by the partition. For non-integral
\(k=\ln(1/\delta)/\ln(1/c)>1\), the detection theorem still gives \(1/k\),
but this partition proof does not establish optimality. For \(k=1\), one
sample suffices in the ideal model; expressions involving
\(\ln(1-1/k)\) are then undefined.

\subsection{Choosing the sample count}

After \(m\) independent draws, the probability of missing invalid work
is at most
\[
 \varepsilon_{\mathrm{sample}}=(1-p)^m.
\]
To make this at most \(2^{-\lambda}\), where \(\lambda>0\) is the
desired error exponent, for \(0<p<1\) choose
\[
 m=\left\lceil
       \frac{\lambda\ln2}{-\ln(1-p)}
    \right\rceil.
\]
Taking logarithms preserves the inequality until division by the
negative number \(\ln(1-p)\), when its direction reverses.

For constant difficulty, checking \(L\) of \(n\) history blocks gives
\(\delta=L/n\). Write \(\log_b x=\ln x/\ln b\) for a logarithm
to base \(b>0\), \(b\ne1\). For fixed \(c\),
\[
 p=\frac{\ln(1/c)}{\ln(n/L)},\qquad
 m=O\bigl(\lambda\log_{1/c}(n/L)\bigr).
\]
The order bound follows from \(-\ln(1-p)\geq p\), which follows by integrating
\(1/(1-t)\geq1\) from zero to \(p\).
An authenticated sample costs one header and an \(O(\log n)\)-node path.
Without sharing duplicate paths, the proof therefore costs
\[
 O\bigl(L+\lambda\log_{1/c}(n/L)\log n\bigr)
\]
header/node units, assuming bounded-size metadata. Actual bytes depend
on header size and metadata encoding, not just this count.

\begin{example}[A checked parameter calculation]\label{ex:fc-samples}
Take \(n=2^{20}\), \(L=2^{10}\), \(c=1/2\), and
\(\lambda=64\). Then \(\delta=2^{-10}\), \(k=10\), and \(p=1/10\).
The least integer sample count is \(m=422\):
\[
 (9/10)^{422}\leq2^{-64}<(9/10)^{421}.
\]
This is an arithmetic example under the stated fork model, not a Zcash
confirmation recommendation. The accompanying numerical check evaluates
both inequalities with exact rational arithmetic.
\end{example}

\section{The interactive verifier}\label{sec:fc-game}

The sampler is useful only when each challenge has a uniquely checkable
answer. We can now state the complete verifier at the abstract
history-interface level.

\begin{construction}[Candidate verification]\label{con:fc-verifier}
The verifier fixes the chain's genesis, header rules, history commitment
scheme, work function, suffix rule, \(c\), and error budget.
\begin{enumerate}
\item Receive a candidate head, the authenticated history length and
  total work, and any enclosing commitment data needed to recover the
  history root. Check canonical encodings and the head's own header
  validity. Freeze this candidate before choosing challenges.
\item Receive and validate the entire required suffix, including its
  ordering, previous-history roots, work and difficulty transitions.
  Check its connection to the head and authenticate its boundary in the
  committed history.
\item Draw \(m\) independent sample points. For each point obtain the
  selected header, its position, sibling metadata, and any adjacent
  header or epoch-boundary data required by the sample-validity rule.
\item Recompute the inclusion path and total weight. Check the work
  coordinate selects that leaf; check the header's own proof of work,
  target and required difficulty data; recover and compare its prefix
  root as in Proposition~\ref{prop:fc-prefix}.
\item Reject the candidate on any failed or missing answer. Otherwise
  retain it for work comparison with the other checked candidates.
\end{enumerate}
\end{construction}

Duplicate samples do not invalidate the bound: the draws remain
independent even if two happen to select one header. A prover may
compress repeated answers, but the verifier must retain the multiplicity
of the original draws when analysing their distribution.

The empty-history and short-history cases use full verification rather
than divide by zero or sample an empty support. A malformed proof is not
a lower-work candidate. It is a rejected candidate. A timeout gives no
positive assertion about chain validity.

\begin{theorem}[Conditional interactive soundness]
\label{thm:fc-interactive}
Assume an honest candidate is available; the underlying chain has stable
prefix security with failure probability
\(\varepsilon_{\mathrm{chain}}\); the authenticated sample interface
fails with probability at most \(\varepsilon_{\mathrm{auth}}\); and
Assumption~\ref{ass:fc-fork} holds with failure probability
\(\varepsilon_{\mathrm{mine}}\).
For one fixed competing candidate outside the short-fork exception, the
probability that the verifier accepts it despite its invalid work is
at most
\[
 \varepsilon_{\mathrm{auth}}
 +\varepsilon_{\mathrm{mine}}
 +(1-p)^m.
\]
The stable-prefix conclusion additionally pays
\(\varepsilon_{\mathrm{chain}}\).
\end{theorem}
\begin{proof}
Exclude the authentication and mining failure events. The frozen
candidate then has well-defined work intervals and enough invalid work.
The suffix either detects it, or Theorem~\ref{thm:fc-detection} gives
detection probability at least \(p\) for each independent challenge.
All \(m\) challenges miss with probability at most \((1-p)^m\).
Adding back the excluded events gives the bound. Stable-prefix agreement
with other honest nodes uses the separate consensus guarantee.
\end{proof}

With several malicious candidates, apply a union bound using their
actual number or a justified bound on it. The server must not be
allowed to reset the experiment indefinitely under the same aggregate
error budget. Availability of an honest proof and verification of
maximal work are parts of the protocol, not consequences of this
single-candidate theorem.

\section{Fiat--Shamir and persistent verification}\label{sec:fc-fs}

Interactive randomness prevents the server from tailoring a committed
history to the sample. A transferable proof instead derives that
randomness from the candidate. This removes interaction, but introduces
the possibility of trying many candidates.

\subsection{Binding the challenge}

The random-oracle model and transcript programming are defined in
\emph{Crypto Guide}, \S\emph{Security contracts and reductions} and
\S\emph{Knowledge, simulation, and Schnorr}. Here the oracle supplies
sampling randomness, not a Schnorr challenge or a witness extractor.

A non-interactive instantiation must hash a domain-separated, canonical
description fixing the network, protocol version, candidate head,
history root, history length, authenticated total work and sampling
parameters. If some field is already unambiguously bound by another,
its duplicate encoding is unnecessary; its binding is not.

The paper derives challenges from the candidate head. Under that
construction, changing the samples entails changing and mining a head.
Nevertheless, the cost of trying heads is a mining-model question.
It is not enough to say that a hash looks random.

\begin{proposition}[A finite grinding budget]\label{prop:fc-grinding}
Suppose each fresh candidate has conditional sample-evasion probability
at most \(\varepsilon\) given the transcript before its fresh oracle
answer. An adversary makes at most \(Q\) such attempts. Its probability
of finding an evading candidate is at most \(Q\varepsilon\).
\end{proposition}
\begin{proof}
For each attempt, its success event has probability at most
\(\varepsilon\), by averaging the conditional bound. The probability of
the union of these events is at most their sum. The candidates need
not be chosen independently.
\end{proof}

Consequently a \(2^{-\lambda}\) aggregate sampling target needs
\[
 (1-p)^m\leq 2^{-\lambda}/Q.
\]
The choice \(Q\) must cover all permitted restarts and candidate-grinding
opportunities. The paper gives a bound in its own mining model; it is
not a universal limit on an implementation's hash queries.

Finite sampling precision creates another explicit error term. Compare
the \emph{selected-header indices}, a finite distribution, with the
indices induced by the ideal continuous draw. A finite set of numerical
coordinates is not close in statistical distance to a continuous
distribution on real coordinates: those supports differ completely.
If each implemented index distribution is within statistical distance
\(\eta\) of its ideal index distribution, independently pair the
implemented and ideal draws so that each pair differs with probability
at most \(\eta\). This pairing is called a coupling; its finite
construction is given in \emph{Math Guide},
\S\emph{Probability and computation}.
The probability that any of the \(m\) pairs differs is at most
\(m\eta\), by the union bound. The security budget must include that
term, or use a separately proved exact discrete index sampler.

\subsection{State retained by a returning client}

A returning client can retain a stable authenticated checkpoint rather
than begin again from genesis. A new candidate must prove consistency
with that checkpoint and satisfy the sampling requirements for its
extension. Old accepted work does not have to be retransmitted merely
because a new suffix arrived.

A current head is not necessarily stable. A client may replace its
unsettled suffix while retaining the earlier common prefix; it must not
turn a previously tentative observation into irrevocable finality.
The appropriate confirmation rule remains the one in
\emph{Consensus Guide}, \S\emph{Reorganisation and confirmation confidence}.

After accepting the header history, the client can request a block's
membership proof and then a transaction's membership proof under that
block's transaction commitment. Effecting-data and authorising-data
commitments must be checked separately where the transaction format
separates them. A transaction-inclusion response alone must not be
treated as evidence of fresh chain selection or transaction execution.

\section{Connecting the construction to Zcash}\label{sec:fc-gap}

Zcash's history commitment is already part of the core consensus
construction. The \emph{Consensus Guide}, \S\emph{Authenticated block
history}, specifies its canonical metadata, its one-block delay,
upgrade-epoch resets, and the outer block commitment. These are inputs
to a possible client verifier, not an implementation of the verifier
above.

The distinction has concrete consequences:
\begin{enumerate}
\item A Zcash history tree covers one upgrade epoch. A proof spanning
  upgrades must authenticate the transitions and recover the relevant
  earlier roots. It cannot interpret the latest root as a genesis-rooted
  tree.
\item The head commits to earlier blocks, not to itself. Its own work,
  validity and total-score contribution must be checked separately.
\item The root uses Zcash's bagging and node encoding. Paper proofs must
  be translated and verified under those exact rules.
\item The root is nested inside the block commitment alongside
  authorising-data commitments. The verifier needs the appropriate
  enclosing values to reconstruct the header's committed digest.
\item Zcash's retarget rule is not the paper's long, clamped-epoch
  model. The aggregate-feasibility predicate and a justified
  valid-fork-weight bound remain necessary work.
\item The wire protocol must specify requests, proofs, integer limits,
  canonical random-bit sampling, error handling, checkpoint persistence
  and multi-server comparison. A node's ability to compute a root does
  not provide those interfaces automatically.
\end{enumerate}

\href{https://zips.z.cash/zip-0221}{ZIP 221}, in its security discussion,
explicitly leaves a stronger analysis for rapidly adjusting chains
before a FlyClient ecosystem is designed. This is a real boundary of
the selected construction, not a historical footnote. The light-client
trust contract examined in \emph{Sync Guide},
\S\emph{Server trust and observable metadata}, does not become trustless
because a block contains a history root.

The result of this volume is therefore precise. The sampling density,
its normalisation, its integer-case optimality and its error accounting
are complete mathematical constructions. Authenticated prefix and weight
checks state the exact interface they consume. Instantiating their
security premises and proof transport for Zcash is not completed by
ZIP 221, by this text, or by the existence of a valid history commitment.

\end{document}
